% MT027 / answer-zh — canonical editable content.
\textbf{随卷答案选B、E。} D错误：Gaussian NB就是用高斯密度对连续特征建模的。


A指课程所讲的标准二分类SVM；要处理多分类，需要另外组合多个分类器，或采用多分类SVM形式。B假设逻辑回归直接使用给定、未经变换的特征向量。C说得过强：特征缩放通常有帮助，尤其是在加入正则化时，但它不是数学上的必要前提。E成立，是因为L1惩罚鼓励稀疏系数。
